\section[Stage 028]{Stage 028: Loaded Axial Profile Selection}
\label{stage:028}

\stagefield{Part / anchor}{Part II; primary anchor \resultanchor{MTDC-T5}.}
\claimstatus{\StatusExactClosure{} inside the two-mode wall profile and rank-one loading closure.}

\stagefield{Purpose}{Stage~028 stops treating the wall-profile angle as an external choice.  The support/mixed branch induces an attractive rank-one load along the D/N overlap vector, and the selected profile is the lower eigenvector of the loaded wall matrix.}

\stagefield{Inputs}{The inputs are the two N/N wall modes, the overlap vector \(v=(\kappa_0,\kappa_1)^T\), and a positive loading coefficient \(\alpha\).}

\stagefield{Verification}{SymPy audit: \StageFile{scripts/moving_throat_pde_stage028_loaded_profile_selection_sympy_audit.py}; transcript: \StageFile{scripts/output/moving_throat_pde_stage028_loaded_profile_selection_sympy_audit.txt}.  Mathematica audit: \StageFile{mathematica/moving_throat_pde_stage028_loaded_profile_selection_mathematica_audit.wl}; transcript: \StageFile{mathematica/output/moving_throat_pde_stage028_loaded_profile_selection_mathematica_audit.txt}.}

\subsection*{Derivation}

\paragraph{1. Bare wall matrix and overlap vector.}
Set
\begin{equation}
\label{eq:app-stage028-qtheta}
q=(q_0,q_1)^T,
\qquad
q_0^2+q_1^2=1,
\qquad
q_0=\cos\theta,
\qquad
q_1=\sin\theta.
\end{equation}
The bare quadrupole wall matrix is
\begin{equation}
\label{eq:app-stage028-Kbare}
K_{\rm bare}=\begin{pmatrix}K_0&0\\0&K_1\end{pmatrix},
\qquad
K_0=K_\eta+6T_\Omega,
\qquad
K_1=K_0+\Delta K_{\rm ax},
\end{equation}
where
\begin{equation}
\Delta K_{\rm ax}=\frac{\pi^2T_w}{L^2}.
\end{equation}
The D/N overlap vector is
\begin{equation}
\label{eq:app-stage028-v}
\boxed{
v=(\kappa_0,\kappa_1)^T,
\qquad
\kappa_0=\frac{2\sqrt2}{\pi},
\qquad
\kappa_1=-\frac{4}{3\pi}}.
\end{equation}

\paragraph{2. Rank-one loaded wall matrix.}
The minimal attractive loading is
\begin{equation}
\label{eq:app-stage028-Keff}
\boxed{
K_{\rm eff}=K_{\rm bare}-\alpha vv^T,
\qquad
\alpha>0}.
\end{equation}
The exact eigenvalues are
\begin{equation}
\label{eq:app-stage028-eigenvalues}
\boxed{
\lambda_\pm=\frac12\left[
K_0+K_1-\alpha(\kappa_0^2+\kappa_1^2)
\pm
\sqrt{\left(\Delta K_{\rm ax}+\alpha(\kappa_0^2-\kappa_1^2)\right)^2+4\alpha^2\kappa_0^2\kappa_1^2}
\right]}.
\end{equation}
The lower selected profile obeys
\begin{equation}
\label{eq:app-stage028-angle-law}
\boxed{
\tan(2\theta_-)
=\frac{2\alpha\kappa_0\kappa_1}{\Delta K_{\rm ax}+\alpha(\kappa_0^2-\kappa_1^2)}}.
\end{equation}
Since \(\kappa_0\kappa_1<0\), positive loading drives \(\theta_-<0\), toward the max-coupling direction and away from the blind-angle branch.

\paragraph{3. Softening threshold.}
The lower eigenvalue reaches zero when \(\det K_{\rm eff}=0\).  This gives
\begin{equation}
\label{eq:app-stage028-alpha-crit}
\boxed{
\alpha_{\rm crit}
=\frac{K_0K_1}{K_1\kappa_0^2+K_0\kappa_1^2}}.
\end{equation}
The stable side of this reduced branch is \(0\leq\alpha<\alpha_{\rm crit}\).

\stagefield{Output}{Stage~028 outputs the rank-one wall matrix \eqref{eq:app-stage028-Keff}, exact eigenvalues \eqref{eq:app-stage028-eigenvalues}, the selected angle law \eqref{eq:app-stage028-angle-law}, and the softening threshold \eqref{eq:app-stage028-alpha-crit}.}

\stagefield{Checks}{\begin{verificationchecklist}
\item The eigenvalues follow from the characteristic polynomial of the \(2\times2\) symmetric matrix \(K_{\rm bare}-\alpha vv^T\).
\item The sign of \(\theta_-\) is fixed because \(\kappa_0>0\) and \(\kappa_1<0\).
\item The threshold is checked by \(\det(K_{\rm bare}-\alpha vv^T)=0\).
\item The loaded branch cannot select the blind angle under positive loading because the selected direction moves toward larger \(|v\cdot q|\), not toward zero overlap.
\end{verificationchecklist}}

\stagefield{Downstream use}{Stage~029 derives \(\alpha\) dynamically from the coupled wall/BdG/Maxwell/mixed Schur complement.  Stages~030--034 use the selected lower eigenvalue and eigenvector overlap to define \(\lambda_-\), \(s_-\), and the softening-depth normal form.}

\clearpage

